%!TEX root = TheoretischePhysikC.tex
% Kern von content/c03-hilbertraum, c03.s8 bis c03.s12. Quellen: docs/plan/script-v0.3/sync-ch4.md

\section{Vorlesung: Operatoren}

\subsection{Operatoren und der Hut}
Das zweite Standbein der Theorie sind die Operatoren: Mit ihnen beschreibt die Quantenmechanik Zustände, Messungen und Dynamik.
\begin{definition}
	Sei $\Hil$ ein Hilbertraum. Ein \emph{Operator} $\op{A}$ auf $\Hil$ ist eine lineare Abbildung $\op{A}\colon\Hil \to \Hil$, das heißt
\begin{equation}
	\op{A}(\ket{\phi} + \ket{\psi}) = \op{A}\ket{\phi} + \op{A}\ket{\psi}, \qquad
	\op{A}(\alpha\ket{\phi}) = \alpha\,\op{A}\ket{\phi}
\end{equation}
für alle $\alpha\in\mathbb{C}$ und alle $\ket{\phi},\ket{\psi}\in\Hil$.
\end{definition}

In dieser Vorlesung ist $\Hil$ endlichdimensional: Es hat eine Basis mit endlicher Indexmenge, wie $\mathbb{C}^n$ mit $\vb{v}\adj\vb{w}$. Mit einer Basis ist ein Operator nichts anderes als eine $n\times n$-Matrix; die Regeln für Matrix mal Spalte sind genau die beiden Bedingungen. Wenn Sie das Wort ,,Operator`` hören, denken Sie an eine Tabelle. Ein Produkt wird von rechts nach links ausgeführt: $\op{A}\op{B}$ heißt erst $\op{B}$, dann $\op{A}$.

Trivial, aber sehr nützlich sind der \emph{Einsoperator} $\mathbb{I}$ mit $\mathbb{I}\ket{\psi} = \ket{\psi}$ und der \emph{Nulloperator} $\widehat{0}$ mit $\widehat{0}\,\ket{\psi} = 0$ für alle $\ket{\psi}\in\Hil$; im endlichdimensionalen Fall sind das die Einheitsmatrix und die Matrix mit lauter Nullen. Einer hingeschriebenen Null sieht man nicht an, was gemeint ist. Wir vereinbaren deshalb:

\begin{merksatz}[Notation]
	Ein \emph{Hut} bedeutet immer einen Operator, $\op{A}$; der \emph{Strich} der Ket-Notation einen Vektor, $\ket{\psi}$; ein Symbol ohne Verzierung eine Zahl. Ohne Hut bleiben der Einsoperator $\mathbb{I}$, die Ket-Bras weiter unten und zwei Sorten von Operatoren, die erst später vorkommen: die Zustände und die Pauli-Matrizen.
\end{merksatz}
Die Vorlesungen ab Vorlesung~5 sind noch in der alten Fassung und schreiben Operatoren oft ohne Hut.

\subsection{Der adjungierte Operator}
Die Kreuz-Operation übertragen wir jetzt auf Operatoren. Der \emph{adjungierte Operator} $\op{A}\adj$ ist durch die folgende Gleichung festgelegt:
\begin{equation}
	\braket{\op{A}\adj\phi}{\psi} = \braket{\phi}{\op{A}\psi} \quad \text{für alle } \ket{\phi},\ket{\psi}\in\Hil.
\end{equation}
Er ist also nicht durch eine Rechenvorschrift bestimmt, sondern durch sein Verhalten unter dem Skalarprodukt. In endlicher Dimension gibt es genau einen solchen Operator. In $\mathbb{C}^n$ mit $\vb{v}\adj\vb{w}$ heißt das Kreuz: erst transponieren, dann alle Einträge komplex konjugieren. Wie der Bra hängt $\op{A}\adj$ aber vom Skalarprodukt ab: Mit einem anderen Skalarprodukt gilt diese Regel für die Einträge in $\mathbb{C}^n$ im Allgemeinen nicht mehr. Die Rechenregeln, die wir ständig verwenden, sind
\begin{equation}
	\begin{gathered}
		(\op{A} + \op{B})\adj = \op{A}\adj + \op{B}\adj, \qquad (z\op{A})\adj = z^*\op{A}\adj, \\
		(\op{A}\op{B})\adj = \op{B}\adj\op{A}\adj, \qquad (\op{A}\adj)\adj = \op{A}.
	\end{gathered}
\end{equation}

\subsection{Klassen von Operatoren}
Mit $\mathbb{I}$ und $\op{A}\adj$ allein definieren wir sechs wichtige Klassen. Ein Operator $\op{A}$ heißt
\begin{enumerate}
	\item \emph{hermitesch}, wenn $\op{A}\adj = \op{A}$,
	\item \emph{positiv}, wenn $\braket{\psi}{\op{A}\psi} \geq 0$ für alle $\ket{\psi}\in\Hil$,
	\item \emph{Projektor}, wenn $\op{A}^2 = \op{A}\adj = \op{A}$,
	\item \emph{isometrisch}, wenn $\op{A}\adj\op{A} = \mathbb{I}$,
	\item \emph{unitär}, wenn $\op{A}\adj\op{A} = \op{A}\op{A}\adj = \mathbb{I}$, und
	\item \emph{normal}, wenn $\op{A}\adj\op{A} = \op{A}\op{A}\adj$.
\end{enumerate}
Jede Klasse hat in der Quantenmechanik ihre eigene Rolle: Die hermiteschen Operatoren beschreiben die Observablen, die positiven, mit einer Normierung, die Zustände, die Projektoren Ja-Nein-Messungen (Ereignis oder kein Ereignis), und die isometrischen und unitären die Dynamik.

Hermitesche Operatoren heißen oft auch selbstadjungiert; in endlicher Dimension ist das dasselbe. Für Operatoren, die nur auf einem Teil des Raumes definiert sind, muss man symmetrisch und selbstadjungiert unterscheiden; das behandelt der Webkurs, ebenso die Isometrien in unendlich vielen Dimensionen, die nicht unitär sind.

$\op{A}$ ist genau dann hermitesch, wenn $\braket{\psi}{\op{A}\psi}$ für jedes $\ket{\psi}$ reell ist; deshalb ist jeder positive Operator hermitesch. Operatoren der Form $\op{B}\adj\op{B}$ sind positiv, weil das Skalarprodukt eines Vektors mit sich selbst nie negativ ist:
\begin{equation}
	\braket{\psi}{\op{B}\adj\op{B}\psi} = \braket{\op{B}\psi}{\op{B}\psi} \geq 0.
\end{equation}
Wegen $\op{P} = \op{P}\adj\op{P}$ ist damit auch jeder Projektor $\op{P}$ positiv. Mit demselben Schritt beginnt der Beweis der Unschärferelation in Vorlesung~9. Ein isometrischer Operator $\op{V}$ lässt das Skalarprodukt invariant, $\braket{\op{V}\phi}{\op{V}\psi} = \braket{\phi}{\psi}$. In endlicher Dimension ist jede Isometrie unitär, und hermitesche, positive und unitäre Operatoren, Projektoren und Isometrien sind normal.

\subsection{Matrixelemente und Ket-Bras}
Sei $\{\ket{\phi_x}\}_{x\in X}$ eine Basis von $\Hil$ mit endlicher Indexmenge. Die \emph{Matrixelemente} von $\op{A}$ in dieser Basis sind die Zahlen $[\op{A}]_{xy}$ in
\begin{equation}
	[\op{A}]_{xy} = \braket{\phi_x}{\op{A}\phi_y}, \qquad
	\op{A}\ket{\Psi} = \sum_{x,y\in X}\ket{\phi_x}\,[\op{A}]_{xy}\,\braket{\phi_y}{\Psi};
\end{equation}
dabei sagt $y$, welcher Basisvektor hineingeht (die Spalte), und $x$, welche Komponente herauskommt (die Zeile). Die zweite Gleichung folgt, wenn man die Entwicklung nach der Basis zweimal einsetzt; über $x$ und $y$ wird unabhängig summiert. Darin kommt keine Matrix vor, nur Basisvektoren und Zahlen: Ein Operator ist bekannt, sobald man seine Matrixelemente kennt. Auch der adjungierte Operator lässt sich ablesen, $[\op{A}\adj]_{xy} = [\op{A}]_{yx}^*$: Für Matrixelemente in einer Orthonormalbasis gilt ,,transponieren und konjugieren`` bei jedem Skalarprodukt.

Kets und Bras kann man selbst als lineare Abbildungen lesen. Ab sofort ist das Ket $\ket{\Psi}$ die Abbildung, die eine Zahl frisst und ein Vielfaches des Vektors $\vb{\Psi}$ produziert, für den es bisher stand. Die Definition des Adjungierten gilt ebenso für lineare Abbildungen zwischen zwei Hilberträumen, hier $\mathbb{C}$ mit dem Skalarprodukt $a^*b$ und $\Hil$. Der adjungierte Operator zum Ket ist der Bra; er frisst einen Vektor und produziert eine Zahl:
\begin{equation}
	\ket{\Psi}\colon \mathbb{C} \to \Hil,\ z \mapsto z\,\vb{\Psi}, \qquad
	\bra{\Psi} = (\ket{\Psi})\adj\colon \Hil \to \mathbb{C},\ \vb{\Phi} \mapsto \braket{\Psi}{\Phi}.
\end{equation}
Bra mal Ket ist dann die Abbildung $z \mapsto \braket{\Phi}{\Psi}\,z$. Wir unterscheiden sie nicht von der Zahl $\braket{\Phi}{\Psi}$; für das Rechnen ändert sich also nichts. In der anderen Reihenfolge, Ket mal Bra, entsteht ein Operator auf $\Hil$, das \emph{Ket-Bra}
\begin{equation}
	\ketbra{\Psi}{\Phi}\,\ket{\chi} = \braket{\Phi}{\chi}\,\ket{\Psi} \quad \text{für alle } \ket{\chi}\in\Hil.
\end{equation}
Ist die Basis aus dem Kontext klar, so trägt der Buchstabe $\phi$ keine Information mehr, und wir schreiben $\ket{x}$ und $\bra{x}$. Damit werden die Entwicklungsformeln zu Operatorgleichungen:
\begin{equation}
	(\ketbra{\Psi}{\Phi})\adj = \ketbra{\Phi}{\Psi}, \qquad
	\mathbb{I} = \sum_{x\in X}\ketbra{x}{x}, \qquad
	\op{A} = \sum_{x,y\in X}\melem{x}{\op{A}}{y}\,\ketbra{x}{y}.
\end{equation}
Prüfen Sie sie, indem Sie beide Seiten auf einen beliebigen Vektor anwenden. Die Gleichung für $\mathbb{I}$ ist eine der hilfreichsten der Quantenmechanik. Physikalisch sagt sie nur, dass Eins gleich Eins ist, aber man darf diese komplizierte Eins an jeder Stelle einer Rechnung einschieben.

\subsection{Drei Operatoren auf \texorpdfstring{$\mathbb{C}^3$}{C3}}
\begin{beispiel}
	Sei $\Hil = \mathbb{C}^3$ mit $\vb{v}\adj\vb{w}$, und seien
	\begin{equation}
		\op{A} = \begin{pmatrix} 0 & 1 & 0 \\ 1 & 0 & 1 \\ 0 & 1 & 0 \end{pmatrix}, \qquad
		\op{B} = \frac13\begin{pmatrix} 1 & 1 & 1 \\ 1 & 1 & 1 \\ 1 & 1 & 1 \end{pmatrix}, \qquad
		\op{C} = \begin{pmatrix} 0 & 1 & 0 \\ 0 & 0 & i \\ -1 & 0 & 0 \end{pmatrix}.
	\end{equation}
	Transponieren und Konjugieren gibt bei $\op{A}$ und $\op{B}$ dieselbe Matrix; beide sind hermitesch. Jeder Eintrag von $\op{B}^2$ ist $\tfrac19(1 + 1 + 1) = \tfrac13$, also ist $\op{B}$ ein Projektor und damit positiv. $\op{A}$ ist dagegen \emph{nicht} positiv: Für $\ket{\psi}$ mit den Einträgen $\tfrac12$, $-\tfrac{1}{\sqrt{2}}$, $\tfrac12$ hat $\op{A}\ket{\psi}$ die Einträge $-\tfrac{1}{\sqrt{2}}$, $1$, $-\tfrac{1}{\sqrt{2}}$, also ist
	\begin{equation}
		\melem{\psi}{\op{A}}{\psi}
		= \begin{pmatrix} \frac12 & -\frac{1}{\sqrt{2}} & \frac12 \end{pmatrix}
		  \begin{pmatrix} - \frac{1}{\sqrt{2}} \\ 1 \\ -\frac{1}{\sqrt{2}} \end{pmatrix}
		= -\sqrt{2}.
	\end{equation}
	$\op{C}$ ist unitär, denn
	\begin{equation}
		\op{C}\adj\op{C}
		= \begin{pmatrix} 0 & 0 & -1 \\ 1 & 0 & 0 \\ 0 & -i & 0 \end{pmatrix}
		  \begin{pmatrix} 0 & 1 & 0 \\ 0 & 0 & i \\ -1 & 0 & 0 \end{pmatrix}
		= \begin{pmatrix} 1 & 0 & 0 \\ 0 & 1 & 0 \\ 0 & 0 & 1 \end{pmatrix},
	\end{equation}
	und ebenso ist $\op{C}\op{C}\adj = \mathbb{I}$. Mit der kanonischen Basis $\ket{0},\ket{1},\ket{2}$ (die Spalten mit der $1$ im ersten, zweiten, dritten Eintrag) ist
	\begin{equation}
		\op{C} = \ketbra{0}{1} + i\,\ketbra{1}{2} - \ketbra{2}{0}.
	\end{equation}
\end{beispiel}

Was die drei Rechnungen gemeinsam haben, sieht man an den Eigenwerten. Ein Eigenwert von $\op{A}$ ist, informell gesagt, eine Zahl $\lambda$, für die $\op{A}$ einen Vektor $\ket{v} \neq 0$ einfach mit $\lambda$ multipliziert, $\op{A}\ket{v} = \lambda\ket{v}$; die genaue Definition kommt in der nächsten Vorlesung. Für jeden Eigenwert, den ein Operator hat, überträgt sich die definierende Gleichung seiner Klasse auf die Zahl $\lambda$:
\begin{equation}
	\begin{aligned}
		&\op{A}\adj = \op{A} &&\Rightarrow\quad \lambda^* = \lambda, \\
		&\braket{\psi}{\op{A}\psi} \geq 0 \text{ für alle } \ket{\psi} &&\Rightarrow\quad \lambda \geq 0, \\
		&\op{A}^2 = \op{A}\adj = \op{A} &&\Rightarrow\quad \lambda^2 = \lambda, \\
		&\op{A}\adj\op{A} = \mathbb{I} &&\Rightarrow\quad \abs{\lambda} = 1.
	\end{aligned}
\end{equation}
Das sind Einschränkungen für die Eigenwerte, die es gibt; die Umkehrung behaupten wir nicht. Die Normalität allein schränkt die Eigenwerte nicht ein. Die Eigenwerte der drei Beispieloperatoren sind
\begin{equation}
	\op{A}\colon\ 0,\ \pm \sqrt{2}, \qquad
	\op{B}\colon\ 1,\ 0,\ 0, \qquad
	\op{C}\colon\ e^{-i\pi/6},\ i,\ e^{7i\pi/6} \quad (\lambda^3 = -i),
\end{equation}
und weitere gibt es nicht.

% fig_vl04_spektren.tex — Operatorklassen an ihren Spektren erkannt:
% wo in der komplexen Ebene die Eigenwerte einer Klasse liegen dürfen,
% und wo die Eigenwerte der drei Beispielmatrizen der Vorlesung landen.
%
% Follows the conventions of fig_vl11_kasten.tex.
%
%   * geometry comes from the defining relations of the lecture, not from
%     a sketch.  For a normal operator the defining equation passes
%     verbatim to each eigenvalue, and that is the whole content of
%     panel (a):
%         A^dag = A          =>  lambda^* = lambda   (the real axis)
%         A^dag A = I        =>  |lambda|  = 1       (the unit circle)
%         A^2 = A^dag = A    =>  lambda^2  = lambda  (the points 0 and 1)
%         <psi|A psi> >= 0   =>  lambda   >= 0       (the half-axis)
%     so the nesting projector < positive < Hermitian is not asserted,
%     it is visible: two points inside a half-line inside a line;
%   * panel (b) plots the spectra of exactly the three matrices of the
%     lecture's Beispiel, from their characteristic polynomials:
%         det(A - lambda) = -lambda(lambda^2 - 2)  =>  0, +-sqrt(2)
%         B = |u><u|, u = (1,1,1)/sqrt(3)          =>  1, 0, 0
%         det(C - lambda) = -(lambda^3 + i)        =>  lambda^3 = -i,
%             i.e. exp(-i pi/6), exp(i pi/2), exp(i 7pi/6) — three points
%             120 degrees apart on the unit circle.
%     The vector psi = (1/2, -1/sqrt2, 1/2) for which the lecture
%     computes <psi|A|psi> = -sqrt(2) is the eigenvector of that
%     eigenvalue, so the number computed in the text is the leftmost blue
%     point: that is why A fails to be positive, and the point is
%     labelled accordingly;
%   * colour encodes the operator class and nothing else — blue
%     Hermitian, green positive, red projector, orange unitary — and it
%     is the same code in both panels, which is what lets the eye carry
%     the map of (a) over to the data of (b);
%   * in (a) the loci ARE the statement, so they are drawn in colour; in
%     (b) they drop back to being reference lines and are drawn neutral
%     grey (`levellinegrey`), with the eigenvalues as the only coloured
%     objects.  Marker shape (disc / ring / diamond) repeats the class
%     distinction, so the plate survives greyscale printing — it also
%     lets the shared eigenvalue 0 of A and B be drawn where it belongs,
%     as a blue disc inside a red ring, instead of being nudged apart;
%   * the classification is a statement about NORMAL operators (the
%     spectral theorem); the grey note in (a) records that normality
%     itself constrains nothing, and the accompanying prose says it.

\begin{center}
\begin{tikzpicture}[x=1.45cm, y=1.45cm,
    % The axes of the complex plane are reference furniture, not data:
    % neutral grey, so the coloured loci of (a) carry the eye.
    cplane axis/.style={-{Stealth[length=4.5pt,width=3pt]},
                        birrengray, line width=0.5pt}]

% ─────────────────────────────────────────────────────────────
%  (a)  The loci: which eigenvalues each class allows
% ─────────────────────────────────────────────────────────────

  % Axes of the complex plane
  \draw[cplane axis] (-2.02,0) -- (2.08,0);
  \draw[cplane axis] (0,-1.50) -- (0,1.88);
  \node[lbl, text=birrengray, font=\scriptsize, anchor=south east]
    at (2.06,0.05) {$\real\lambda$};
  \node[lbl, text=birrengray, font=\scriptsize, anchor=west]
    at (0.09,1.78) {$\imag\lambda$};

  % positive:  lambda >= 0 — a highlighter stroke under the Hermitian line
  \draw[line width=6pt, birrengreen!55, line cap=butt] (0,0) -- (1.80,0);
  % Hermitian: lambda^* = lambda
  \draw[line width=1.4pt, cgblue] (-1.80,0) -- (1.80,0);
  % unitary:   |lambda| = 1
  \draw[line width=1.3pt, birrenorange] (0,0) circle (1);
  % projector: lambda^2 = lambda
  \fill[birrenred] (0,0) circle (3pt);
  \fill[birrenred] (1,0) circle (3pt);
  \node[lbl, text=birrenred, font=\scriptsize, anchor=north east] at (-0.04,-0.04) {$0$};
  \node[lbl, text=birrenred, font=\scriptsize, anchor=south west] at ( 1.05, 0.04) {$1$};

  % Class labels, each in its own colour, each outside the unit circle
  \node[lbl, text=birrenorange, font=\scriptsize, align=center, anchor=south]
    at (0,1.07) {\tpcifenglish{unitary}{unitär}\\$|\lambda|=1$};
  \node[lbl, text=birrenred, font=\scriptsize, align=right, anchor=south east]
    at (-1.04,0.12) {\tpcifenglish{projector}{Projektor}\\$\lambda\in\{0,1\}$};
  \node[lbl, text=cgblue, font=\scriptsize, align=right, anchor=north east]
    at (-1.04,-0.10) {\tpcifenglish{Hermitian}{hermitesch}\\$\lambda^{*}=\lambda$};
  \node[lbl, text=birrengreen, font=\scriptsize, align=left, anchor=north west]
    at (1.04,-0.10) {\tpcifenglish{positive}{positiv}\\$\lambda\ge 0$};
  \node[lbl, text=birrengray, font=\scriptsize, anchor=north west]
    at (0.12,-1.14) {normal: \tpcifenglish{all of}{ganz} $\C$};

  \node[lbl head, anchor=south] at (0,2.02)
    {(a) \tpcifenglish{where the eigenvalues lie}{wo die Eigenwerte liegen}};

% ─────────────────────────────────────────────────────────────
%  (b)  The same plane, with the spectra of the lecture's example
% ─────────────────────────────────────────────────────────────
\begin{scope}[xshift=6.95cm]

  % Reference loci, now neutral: the unit circle and the real axis
  \draw[levellinegrey] (0,0) circle (1);
  \draw[cplane axis] (-2.02,0) -- (2.08,0);
  \draw[cplane axis] (0,-1.50) -- (0,1.88);
  \node[lbl, text=birrengray, font=\scriptsize, anchor=south east]
    at (2.06,0.05) {$\real\lambda$};
  \node[lbl, text=birrengray, font=\scriptsize, anchor=west]
    at (0.09,1.78) {$\imag\lambda$};
  \node[lbl, text=birrengray, font=\scriptsize, anchor=north west]
    at (0.10,-1.06) {$|\lambda|=1$};

  % C:  lambda^3 = -i  — three points 120 degrees apart on the circle
  \foreach \a in {-30,90,210}{
    \fill[birrenorange] (\a:1) ++(0,3.4pt) -- ++(3.4pt,-3.4pt)
                        -- ++(-3.4pt,-3.4pt) -- ++(-3.4pt,3.4pt) -- cycle;
  }
  % A:  lambda = 0, +-sqrt(2)
  \fill[cgblue] ({-sqrt(2)},0) circle (2.5pt);
  \fill[cgblue] (0,0)          circle (2.5pt);
  \fill[cgblue] ({ sqrt(2)},0) circle (2.5pt);
  % B:  lambda = 1, 0, 0  — the ring at 0 encircles A's disc
  \draw[birrenred, line width=1.1pt] (0,0) circle (4.3pt);
  \draw[birrenred, line width=1.1pt] (1,0) circle (4.3pt);

  % Eigenvalue readings
  \node[lbl, font=\scriptsize, anchor=north]      at ({-sqrt(2)},-0.09) {$-\sqrt{2}$};
  \node[lbl, font=\scriptsize, anchor=north]      at ({ sqrt(2)},-0.09) {$\sqrt{2}$};
  \node[lbl, font=\scriptsize, anchor=north]      at (1,-0.19)          {$1$};
  \node[lbl, font=\scriptsize, anchor=north east] at (-0.05,-0.12)      {$0$};

  % Operator labels, in the colour of the class each one belongs to
  \node[lbl, text=cgblue, font=\scriptsize, align=right, anchor=south east]
    at (-1.06,0.44) {$\widehat A$\\\tpcifenglish{Hermitian}{hermitesch}};
  \draw[cgblue, line width=0.4pt] (-1.45,0.42) -- (-1.42,0.08);
  \node[lbl, text=birrenorange, font=\scriptsize, align=center, anchor=south]
    at (0,1.10) {$\widehat C$\\\tpcifenglish{unitary}{unitär}};
  \node[lbl, text=birrenred, font=\scriptsize, align=left, anchor=south west]
    at (1.06,0.34) {$\widehat B$\\\tpcifenglish{projector}{Projektor}};
  \draw[birrenred, line width=0.4pt] (1.14,0.32) -- (1.05,0.09);

  % Why the lecture bothers to compute <psi|A|psi> = -sqrt(2)
  \node[lbl, text=birrengray, font=\scriptsize, anchor=north]
    at ({-sqrt(2)},-0.70) {\tpcifenglish{not positive}{nicht positiv}};

  \node[lbl head, anchor=south] at (0,2.02)
    {(b) \tpcifenglish{the example's $\widehat A,\widehat B,\widehat C$}%
                      {$\widehat A$, $\widehat B$, $\widehat C$ aus dem Beispiel}};

\end{scope}
\end{tikzpicture}
\end{center}


Teil (a) der Abbildung zeigt, wo die Eigenwerte einer Klasse liegen dürfen, Teil (b) die Eigenwerte von $\op{A}$, $\op{B}$ und $\op{C}$. Dort liegt $-\sqrt{2}$ links der Null; genau daran scheitert die Positivität von $\op{A}$, denn $\ket{\psi}$ von oben ist ein Eigenvektor der Norm eins zu diesem Eigenwert.
